\scanpage{148}
\mspage{75}
% archive sheet number: 147
% L1–6
que, ou bien les $\chi_i$ sont égaux -- et \inserted{alors} égaux \deleted{\unsure{au}}\inserted{à l'unique} caractère de $X$ laissé fixe par l'action précédente -- ou bien $\chi_i$ prend deux valeurs $\chi,\chi'$, correspondant à \deleted{\unsure{deux}} la partition de $\Pi$ en deux classes suivant le \deleted{\unsure{caractè}} ss-groupe de $Z'_n\simeq Z$ noyau du caractère
% L7
\[
\varepsilon: Z'_n\longrightarrow\{\pm1\}
\]
% L8–15
qui décrit l'action de $Z'_n$ sur $X$. Mais il faudrait déjà se rendre compte s'il peut bien arriver que l'action de $Z'_n$ sur $X$ soit non triviale -- dans le cas congruentiel défini par \deleted{\unsure{$G_0$}} un module libre de rg 2 $M$ sur $\Lambda_n$, le premier cas où la question se pose est pour $n=\deleted{\unsure{8|3|5}}\unsure{5|8}$, $Z'_n\simeq\deleted{\unsure{2}}\pm1$, quelle est son action sur $[\mathrm{Gl}'(2,\mathbb Z/\deleted{\unsure{8|5}}\unsure{5|8}\mathbb Z)^{\pm1}]_{\rm ab}$ ? Tant qu'à faire, autant passer à la limite sur $n$ et se demander : quelle est, \deleted{l'action de relèvement}\inserted{pour la structure d'extension}
% L18
\[
1\longrightarrow \mathrm{Gl}'(2,\widehat{\mathbb Z})^{\pm}
 \longrightarrow\mathrm{Gl}'(2,\widehat{\mathbb Z})
 \mathrel{\unsure{-|\longrightarrow}}\widehat{\mathbb Z}^{*}/\pm1\mathrel{-}1
\tag*{(52)}
\]
% L19–21
l'\deleted{\unsure{structure d'extension}}\inserted{opération correspondante}
\[
\left\{\begin{aligned}
&\text{opération de }\widehat{\mathbb Z}^{*}/\{\pm1\}\text{ sur }
[\mathrm{Gl}'(2,\widehat{\mathbb Z})^{\pm}]_{\rm ab}
 \simeq\underbrace{\mathbb F_2\times\mathbb F_2}_{\text{base }\dot\tau_1,\dot\tau_\infty}\\
&\text{i.e. l'homomorphisme}\\
&\hspace{1em}\widehat{\mathbb Z}^{*}/\{\pm1\}\longrightarrow
\unsure{\mathfrak S|\mathscr S}_{\{\dot\tau_1,\dot\tau_\infty,
 \underbrace{\dot\tau_1\dot\tau_\infty}_{\substack{\deleted{\unsure{\dot\rho_0=\dot\tau_0}}\\
 \unsure{\dot\lambda_0^{-1}=\dot\varepsilon_0|\dot\lambda_0=\dot\varepsilon_0}}}\}}\simeq\mathfrak S_3
\end{aligned}\right.
\tag*{(53)}
\]
% L22–25
\deleted{\unsure{j}}il est clair, en regardant le sous-groupe de $G(2,\widehat{\mathbb Z})$ formé des matrices
\[
\left\{\begin{pmatrix}\lambda&\mu\\0&1\end{pmatrix}\ \bigg|\ 
\lambda\in\widehat{\mathbb Z}^{*},\ \mu\in\widehat{\mathbb Z}\ ,\right.
\]
qui est produit $1/2$-direct de $\widehat{\mathbb Z}^{*}$ \deleted{\unsure{de}} (isom au groupe des $\left(\begin{smallmatrix}\lambda&0\\0&1\end{smallmatrix}\right)$
